%% kapa: exam=2 round=36 year=2025 period=2 subject=appraisal_theory subject_ko=감정평가이론 seq=2 number=2 points=30
%% source: Q-Net 감정평가사 기출문제 (cst00309) artlSeq=5239325 file=raw/1008/5239325/2025년 제36회 감정평가사 2차 감정평가이론.pdf
%% conversion: kordoc 4.24.51 -> tools/gen2.py (auto); review=proofread-gpt-sol; flags=none
\begin{longquestion}{【문제 2】}
“감정평가”란 토지등의 경제적 가치를 판정하여 그 결과를 가액(價額)으로 표시하는 것을 말한다. 다음 물음에 답하시오. (30점) \par
\medskip\noindent 물음 1) 감정평가의 정의(定義)에 규정된 ‘경제적 가치’와 ‘판정’의 의의를 각각 쓰시오. (10점) \par
\medskip\noindent 물음 2) 관련 법령에 따라 ‘복수(複數) 감정평가’를 하는 공시지가평가, 보상평가, 재개발사업 종전자산평가의 경제적 가치 판정에 있어 유의할 점을 각각 쓰시오. (단, 감정평가 대상물건은 토지에 한함) (20점) \par
\begin{framedtext}
\centerline{<참고 사항>}
○ 공시지가평가: ｢부동산 가격공시에 관한 법률｣에 따른 표준지공시지가 조사ㆍ평가 \par
○ 보상평가: ｢공익사업을 위한 토지 등의 취득 및 보상에 관한 법률｣등 법령에 따라 공익사업을 목적으로 취득하는 토지에 대한 손실보상을 위한 감정평가 \par
○ 재개발사업 종전자산평가: ｢도시 및 주거환경정비법｣에 따른 재개발사업의 관리처분계획 수립과 관련된 종전자산의 감정평가 \par
\end{framedtext}
\end{longquestion}
